\documentclass{ximera}

\title{Approximating Areas under Curves}
\author{MATH 425: Calculus I}

\begin{document}
\begin{abstract}
   Working with the peers in your group, solve the following problems. Make sure to show and justify all your work. Make sure everyone in the group understands the solution and participates. Be prepared to report your answers to the whole class. 
\end{abstract}
\maketitle


\begin{exercise}
    Consider the function $f(x)=3\sin(x)$ on the interval from $[0,\pi]$.  Estimate the area between $f(x)$ and the $x$-axis by \textit{formulating} $R_4$ and $L_4$.\\
\end{exercise}

\begin{exercise}
    Suppose that an object moving along a straight line path has its velocity in feet per second at time 
 in seconds given by $v(t)=\frac29(t-3)^2+2$.
 \begin{enumerate}
    \item Sketch the region whose exact area will tell you the distance the object travelled over the time interval $2\leq t\leq 5$.
    \item Estimate the distance travelled on $[2,5]$ by computing $L_3$, $M_3$, and $R_3$.
    \item Does averaging $L_3$ and $R_3$ result in the same value as $M_3$?  If not, what do you think the average of $L_3$ and $R_3$ measures?
    \item or this question, think about an arbitrary function $f$, rather than the particular function $V$ given above. If $f$ is positive and increasing on $[a,b]$, will $L_n$ over-estimate or under-estimate the exact area under $f$ on $[a,b]$?  Will  $R_n$ over- or under-estimate the exact area under $f$ on $[a,b]$?Explain.
 \end{enumerate}
 Problem adapted from Boelkins, et al., \textit{Active Calculus}, 2nd ed.
 

\end{exercise}

\begin{exercise}
    Use the following formulas to calculate the exact area under $g(x)=x^2-4x$ from $x=1$ to $x=3$.  $$\sum_{j=1}^{n} j=\frac{n(n+1)}{2}$$
    $$\sum_{j=1}^n j^2=\frac{n(n+1)(2n+1)}{6}$$
\end{exercise}   

\end{document}
